\documentclass[a4j]{jarticle}
\usepackage{ascmac}
\usepackage[dvipdfmx]{graphicx}

\begin{document}

\vfill
\begin{screen}
\begin{center}
\huge\bf 代数第一回レポート課題\\
\vspace{5mm}
\large\bf
学籍番号　2022311042 \quad 名前 神野友哉\\
\today
\end{center}
\end{screen}
\thispagestyle{empty}

%
\section{イミテーションゲーム}
\subsection{この映画を選んだ理由}
三つの挙げられていた映画を調べ、暗号の解読に成功したという
文言に興味がわいたから。

\subsection{この映画の全体像}
イミテーションゲームは2014年にイギリスで作成された映画であり、
第二次世界大戦中に枢軸国の一国であるドイツがメッセージ伝達手段として
用いたエニグマ暗号を解読することに成功した実在の人物、アランチューリング
に焦点を当てて描画している。彼の過去と、現代を織り交ぜながら描写している
ことが、作品において特徴的である。
\newpage
\subsection{あらすじ}
映画が始まって最初は、1951年のイギリスの都市マンチェスターにおいて、アランチューリングの家に泥棒が
入ったという情報の元、刑事がその家に向かうと、その家の主である彼に追い出されます。
刑事は当然、彼がやましい事をしているのではないかと訝しみ、調査が進むに連れ、
彼は当時のイギリスで違法とされていた、同性愛者であることが判明します。更に、ソビエト連邦のスパイである
疑惑も掛けれれてしまいました。彼は、刑事に尋問を受ける中、過去を語り始めました。
その後、物語は1939年の第二次世界大戦中のイギリスで始まります。彼は若い数学者です。
彼は、採用面接を受けるため、ブレッチリー無線機器製造所を訪ね、海軍のデニストン中佐に面接を受けます。
中佐は彼のことを、不適合だと考えて採用したくはありませんでしたが、彼はその採用の意味が
エニグマ解読機の作成であることを推測しており、それを言ってしまいました。そうして、彼を採用せざる得なくなり、
秘密裏にブレッチリー・パークと呼ばれる政府に隠匿された場所で、ナチス・ドイツが
使用する暗号機「エニグマ」の解読に取り組み始めます。彼以外の他のメンバーは、
エニグマ暗号を骨抜きにするため、パターンを探すことに着手しますが、彼は、エニグマの暗号を解読するための
「マシン」を開発することに着手します。なぜなら、エニグマが生成する数千もの暗号パターンを解読するために、
人間ではなく機械を用いるしかないと考えたからです。
ここで、一旦時代が1951年に移ります。刑事が彼アランチューリングのことを調べても調べても、
彼の情報が調査されなかった主な理由は、彼が第二次世界大戦中にブレッチリー・パークで
暗号解読の仕事を行っていたため、その活動が極秘扱いされていたからです。
ブレッチリー・パークはイギリス政府の秘密機関であり、エニグマ暗号の解読に関わる活動が行われていました。
このような秘密情報や暗号解読の活動は国家の安全保障に関わる重要なものであり、一般の刑事調査の範囲外にありました。
これが、刑事の彼に対する不信感を増長させてしまった主な原因です。
また、舞台が戻り1939年になり、彼は独自の暗号解読機を作り上げることを目指す
一方、彼は社交性に欠け、対人関係に苦手意識を持っていました。
彼のコミュニケーションスタイルは他のメンバーと衝突することがあり、
特にヒュー・アレクサンダーとは意見が対立しました。彼は当時の英国首相ウィンストンチャーチルに手紙を送ることで、
自らにチームを率いる権限を得ることに成功し、チームの内二名を止めさせてしまいました。
その後、クロスワードパズルで優秀な人を採用しようと試み、その際、唯一の女性であった、
ジョーンクラークを採用します。彼は、チームの他のメンバーとはジョーンクラークの助言もあって
関係はある程度改善します。その後、デニストン中佐が解読機の成果が芳しくないことを理由に、彼をチームから除外
しようとしましたが、もし彼が止めたとしたらチームの他のメンバーも止めると宣言し、猶予期間を得ることに成功します。
紆余曲折ありましたが、チューリングのチームはマシンの完成を目指し続けて、ついにエニグマの解読に成功します。
これにより、連合軍はドイツの通信を傍受し、戦局に大きな影響を与えることができました。
一方、1951年では、スパイの容疑は晴れましたが、彼は当時違法であった同性愛のおおよそ全てを刑事に話してしまい、
わいせつ罪で起訴され、２年の懲役か薬剤による去勢を選ぶことを余儀なくされ、
後者を選択しました。その後、それの影響もあってか彼は若くして自殺し、その生涯を終えました。
\newpage
\subsection{感想・コメント}
実話と聞いて、とても驚きましたが、全て実話だと逆に失望するような人格だったと思いました。
だからこそ、映画の中のオリジナル要素として、チャーチルに手紙を送り、
リーダーに任命されることで、チームの二名を追放するという内容は、フィクションで良かったと思いました。
途中から、交互に現代と過去を書くのが面倒になって、多少適当に書いてしまいました。さらに、主人公のアランチューリングの
年少期も描写されていますが、十分人間関係が苦手という描写があると感じ省きました。暗号との出会いもこの内容の
中に含まれていることですが。最後に、映画を見た後に内容と順序をかなり忘れてしまい、ネットのあらすじまとめを参考にしました。
映画の中身をレポートに起こすということが、如何に難しいかを体験したと思います。それこそ、他者との内容のかぶりを排除しようと
試みても、絶対に内容を被らせないといったことは不可能であり、また、何を書けばいいんだ？と常々考えていました。
そのうえ、小説等の文字データとは異なり、時間でコマが進む様な分類のデータは、再度巻き戻って見返す際に時間リソースを
大きく割いてしまう為に、不便極まりないと結論付けました。

\subsection{参考サイト}
https://mihocinema.com/imitation-game-8141
映画『イミテーション・ゲーム エニグマと天才数学者の秘密』のネタバレあらすじ結末と感想

\end{document}